{\small
\rowcolors{2}{rowlt}{white}
\begin{longtable}{>{\raggedleft}p{5mm} p{12mm} >{\raggedright}p{85mm} r r r p{14mm}}
\toprule
\rowcolor{hdrblue}\hdr{It} & \hdr{Type} & \hdr{Change and reason} & \hdr{Cycle} & \hdr{$\Delta$s} & \hdr{Viol.} & \hdr{Result}\\
\midrule
\endhead
1 & tool & \textbf{slow vertical 50 mm approach/retreat at every grip/release (tool)}\newline Why: 12 of v3's 18 approach violations (r3) were slow vertical legs shorter than 50 mm; one tool fixes all of that kind at once. Accepted +6.8 s, because legal comes before fast. & 246.6 & +6.8 & 17 & kept \\
2 & mixed & \textbf{guide gripped from above - grip pocket becomes a vertical through-notch (mixed: jig\_cad + program)}\newline Why: 4 of the remaining 6 r3 hits were the guide's sideways slide into its closed side pocket. A vertical through-notch allows a grip from above (DFA: top-down handling). CAD + program together, because neither half works alone. & 247.5 & +0.9 & 8 & kept \\
3 & layout & \textbf{plank jig moved 60 mm toward base A (layout + plank-pick moves shifted)}\newline Why: 5 of 8 violations were the forearm (link 3) hitting planks 4/5 in the jig while it reached for the guide on base B. Moving the jig 60 mm away from that reach removes the cause instead of re-routing around it. & 247.0 & -0.5 & 4 & kept \\
4 & program & \textbf{frame lift-out PTP (Point82) $\to$ fast Line}\newline Why: r4 flagged the PTP lift-out starting only 50 mm above the frame grip. Both ends have the same tool orientation, so a straight Line rises vertically and is legal. & 247.9 & +0.9 & 3 & kept \\
5 & mixed & \textbf{raise the plank1 transit (Point56) so plank1 clears plank2}\newline Why: the carried plank1 swept through the top of the already placed plank2 (r1). Measured it first (L07): needed more transit height, not a new route. & 247.9 & +0.0 & 2 & kept \\
6 & program & \textbf{vertical exit out of the frame release}\newline Why: the exit after the frame release was a slow sideways dogleg ($\sim$18 s at 20 mm/s, and an r3 hit). Leaving straight up the way it came in removes both. & 229.6 & -18.3 & 1 & kept \\
7 & jig\_cad & \textbf{relief recesses under the plank ends in base B}\newline Why: the last violation was a zero-depth contact of the plank ends on the base B pocket floor. A 2 mm relief under the plank ends should remove it. Rejected: it also removed the floor under 4 interface pieces (r6). & 229.6 & +0.0 & 4 & rejected \\
8 & jig\_cad & \textbf{plank-end relief in base B that keeps the floor under the pieces B}\newline Why: iter 07's idea was right but the cut was too wide. Cut only the plank-end footprints minus the interface-piece footprints. First legal design. & 229.6 & +0.0 & 0 & kept \\
9 & program & \textbf{drop the redundant plank push after every plank release}\newline Why: biggest phase = planks (127 s). The push after each release (close, push 20 mm, open) moved nothing: the release pose was already home (r8 still passes). Remove a whole handling step x6. & 217.7 & -11.9 & 0 & kept \\
10 & program & \textbf{shorten the slow plank-place descent from 75 to 50.5 mm}\newline Why: the place descent was slow for 75 mm but the rule needs only the last 50 mm slow; 25 mm at 20 mm/s $\to$ at 100 mm/s saves $\sim$1 s per plank. & 211.8 & -5.9 & 0 & kept \\
11 & program & \textbf{outbound guide transfer with PTPs instead of long 100 mm/s Lines}\newline Why: guide phase (62 s) was long 100 mm/s Lines; a PTP over the same distance is faster (joint speed). Outbound transfer first. & 202.8 & -9.0 & 0 & kept \\
12 & program & \textbf{return guide transfer with PTPs instead of long 100 mm/s Lines}\newline Why: same evidence and change as iter 11, for the return transfer ($\sim$22 s of Lines). & 194.6 & -8.2 & 0 & kept \\
13 & program & \textbf{lower the plank transit points from TCP z 391 to grip/release + 102.5 mm}\newline Why: every plank pick/place went through a transit at TCP z 391 although the rules only need $\sim$100 mm clearance; lower transits shorten 24 fast Lines. & 190.8 & -3.8 & 0 & kept \\
14 & program & \textbf{frame-1 approach with one PTP + vertical Line instead of the Lines via (-500, 0)}\newline Why: the frame-1 approach was 11.5 s of Lines and small PTPs via a detour; one PTP + one Line covers it. & 186.6 & -4.2 & 0 & kept \\
15 & program & \textbf{retreat after the frame release as Line up + one PTP}\newline Why: the retreat after the frame release was 3 Lines + a PTP (5.7 s); Line up + one PTP is shorter. & 185.5 & -1.1 & 0 & kept \\
16 & program & \textbf{plank pick slow lift 60 $\to$ 50.5 mm}\newline Why: still biggest phase planks; the slow lift after each plank grip was 60 mm where 50.5 mm is enough (same logic as iter 10, on the pick side). & 183.2 & -2.3 & 0 & kept \\
17 & program & \textbf{frame release slow leg 75 $\to$ 50.5 mm}\newline Why: same 75 $\to$ 50.5 mm slow-leg cut for the frame release. & 182.2 & -1.0 & 0 & kept \\
18 & program & \textbf{lift after the guide release on base A as a short vertical Line + the existing PTP}\newline Why: after releasing the guide the arm climbed to z 461, a height left over from v3's route; climb only as high as the next PTP needs. & 181.2 & -1.0 & 0 & kept \\
19 & program & \textbf{J6 unwrapped +360 deg from the guide flip to the plank6 pick and over the guide return}\newline Why: the longest PTP was J6-bound (232 deg); unwrapping J6 by +360 deg keeps the same pose with a shorter turn (lowering a PTP was tried first in screens and dropped, L19). & 176.9 & -4.3 & 0 & kept \\
20 & program & \textbf{guide return grip at J6 +180 (short way from plank4), lift Point11r higher so the opposite flip clears}\newline Why: the guide return grip turned J6 the long way from plank4; grip at J6 +180 and lift higher so the opposite flip clears plank6. & 175.4 & -1.5 & 0 & kept \\
21 & program & \textbf{guide placement on base A: PTP end Point11h offset sideways and lower, short diagonal Line down}\newline Why: a 244 mm Line dropped the guide straight down only because the PTP end had to stay 100 mm from a plank release point; move that PTP end sideways instead of up. & 175.0 & -0.4 & 0 & kept \\
22 & program & \textbf{plank transits offset sideways out of the plank's thin xy box, lowered to event + 102}\newline Why: plank transits sat at +117.5 mm because they were inside the plank's thin (5 mm) xy box; a 10 mm side offset leaves the box and lets them drop to +102 mm (24 Lines shorter). & 171.7 & -3.3 & 0 & kept \\
23 & program & \textbf{guide-grip transits over base A (Point15v, Point14h) offset sideways and lowered}\newline Why: same side-offset idea for the guide-grip transits over base A (long 240 mm vertical Lines). & 170.0 & -1.7 & 0 & kept \\
24 & program & \textbf{home pose moved onto the joint path between the last release lift and the first guide approach}\newline Why: start/end of the cycle cost 4 s of wrist-branch (J5) swings; put home on the joint path between the last lift and the first approach. & 165.8 & -4.2 & 0 & kept \\
25 & program & \textbf{the guide stays on base B (no guide trip to base A and back)}\newline Why: the guide's trip to base A and back ($\sim$46 s) only serves as a placing aid; frame 1 already holds the planks upright alone, and no rule needs the trip. Lesson L25: question whether a phase is needed before tuning it. & 118.3 & -47.5 & 0 & kept \\
26 & mixed & \textbf{plank jig moved toward base A (shorter pick $<$$\to$ place PTPs)}\newline Why: with the guide trip gone, the gap between plank jig and base A was free; moving the jig closer shortens 12 pick/place PTPs (got less than expected: those PTPs were J6-bound). & 117.5 & -0.8 & 0 & kept \\
27 & mixed & \textbf{planks 4-6 presented pre-rotated 80 deg in the jig, placed as a group after 3, 1, 2}\newline Why: five PTPs were J6-bound because planks 4-6 sit 80 deg turned in the frame vs the jig; present them pre-rotated (DFA: present parts in assembly orientation) and place them as a group. & 116.2 & -1.3 & 0 & kept \\
28 & program & \textbf{frame flip as one PTP onto the other wrist branch (short return home)}\newline Why: the end of the cycle was set by a 228 deg J5 swing; a PTP flip onto the other wrist branch (a Line could not, L24) shortens the return home. & 115.8 & -0.4 & 0 & kept \\
29 & program & \textbf{frame-grip approach PTP ends further down the diagonal Line (short Line)}\newline Why: the 320 mm diagonal Line to the frame-foot grip was the longest single move left; end the PTP further down it. & 113.4 & -2.4 & 0 & kept \\
30 & program & \textbf{shorter fast lift Point80u $\to$ Point80L after the frame-foot grip}\newline Why: the lift after the frame grip climbed higher than r4 needs before the flip PTP; stop it lower (small, -0.5 s). & 112.9 & -0.5 & 0 & kept \\
\bottomrule
\end{longtable}
}
